\documentclass[12pt]{article}
\usepackage[T1]{fontenc}
\usepackage[utf8]{inputenc}
\usepackage{lmodern}
\usepackage{textcomp}
\usepackage{enumitem}
\usepackage{geometry}
\geometry{margin=1in}
\begin{document}
\begin{center}
Answer Key - Biology - Undergraduate Level Assessment 4 \
\small (Answers are ordered exactly as in the question paper.)
\end{center}

\section*{Multiple Choice}
\begin{enumerate}[label=\arabic*.]
\item \textbf{Answer:} A
\\ \textit{Options:} A) DNA migrates from positive charge to negative charge.; B) Smaller DNA travels faster.; C) The DNA migrates only when the current is running.; D) The longer the current is running, the farther the DNA will travel.
\\ \textit{Note:} Correct option: A - DNA migrates from positive charge to negative charge.
\item \textbf{Answer:} A
\\ \textit{Options:} A) hydrogen; B) methane; C) nitrogen; D) sulfur dioxide; E) neon
\\ \textit{Note:} Correct option: A - hydrogen
\item \textbf{Answer:} A
\\ \textit{Options:} A) Complete absence of segmentation; B) Hindrance in the development of the nervous system; C) Inability to undergo metamorphosis; D) Absence of a group of contiguous segments; E) Tumor formation in imaginal discs
\\ \textit{Note:} Correct option: A - Complete absence of segmentation
\item \textbf{Answer:} B
\\ \textit{Options:} A) cross section of muscle tissue; B) longitudinal section of a shoot tip; C) longitudinal section of a leaf vein; D) cross section of a fruit; E) cross section of a leaf
\\ \textit{Note:} Correct option: B - longitudinal section of a shoot tip
\item \textbf{Answer:} B
\\ \textit{Options:} A) Vinegar has a lower concentration of hydrogen ions than hydroxide ions; B) Vinegar has an excess of hydrogen ions with respect to hydroxide ions; C) Vinegar has a higher concentration of hydroxide ions than water; D) Vinegar has no hydroxide ions; E) Vinegar does not contain hydrogen ions
\\ \textit{Note:} Correct option: B - Vinegar has an excess of hydrogen ions with respect to hydroxide ions
\end{enumerate}

\section*{True / False}
\begin{enumerate}[label=\arabic*.]
\setcounter{enumi}{5}
\item \textbf{Answer:} False
\\ \textit{Options:} A) True; B) False
\\ \textit{Note:} LLM-generated false statement. Correct answer: 'decreases as blood pH decreases'.
\item \textbf{Answer:} False
\\ \textit{Options:} A) True; B) False
\\ \textit{Note:} LLM-generated false statement. Correct answer: 'The thick muscular wall of the stomach'.
\item \textbf{Answer:} True
\\ \textit{Options:} A) True; B) False
\\ \textit{Note:} LLM-generated true statement based on MCQ.
\item \textbf{Answer:} True
\\ \textit{Options:} A) True; B) False
\\ \textit{Note:} LLM-generated true statement based on MCQ.
\item \textbf{Answer:} False
\\ \textit{Options:} A) True; B) False
\\ \textit{Note:} LLM-generated false statement. Correct answer: 'Silk worms spin cocoons based on the phase of the moon'.
\end{enumerate}

\section*{Long Form Response}
\begin{enumerate}[label=\arabic*.]
\setcounter{enumi}{10}
\item \textbf{Answer:} stabilizing selection. Explain why this is the correct answer and provide supporting evidence. Reference key facts or concepts that support this choice.
\\ \textit{Note:} Long-form derived from MCQ; award credit for stating the correct idea and giving supporting reasoning.
\item \textbf{Answer:} A flock of birds can be considered an elementary form of society called a motion group. They communicate with one another in order to stay together as they move from place to place. The cooperative behavior enables them to more efficiently detect and avoid predators.. Explain why this is the correct answer and provide supporting evidence. Reference key facts or concepts that support this choice.
\\ \textit{Note:} Long-form derived from MCQ; award credit for stating the correct idea and giving supporting reasoning.
\end{enumerate}

\vfill
\noindent\textit{End of Answer Key}
\end{document}